\chapter{Введение}

\section{Цель}  

Произвести настройку маршрутизаторов. Научиться
настраивать протокол RIP на маршрутизаторах для организации сетевой
связности.
\section{Задачи}

\begin{enumerate}
\item интерфейс Gi1/0/1 маршрутизатора R1 соединить с интерфейсом Gi1/0/2
маршрутизатора R2.
\item Интерфейс Gi1/0/2 маршрутизатора R1 соединить с интерфейсом Gi1/0/1
маршрутизатора R3.
\item Интерфейс Gi1/0/1 маршрутизатора R2 соединить с интерфейсом Gi1/0/2
маршрутизатора R3.
\item Компьютер VPC1 подключить к интерфейсу Gi1/0/3 маршрутизатора R1.
\item Компьютер VPC2 подключить к интерфейсу Gi1/0/3 маршрутизатора R2.
\item Компьютер VPC3 подключить к интерфейсу Gi1/0/3 маршрутизатора R3.
\item Настроить имена устройств и IP - адреса интерфейсов согласно таблице
5.2.1.
\item Настроить протокол RIP и включить анонсирование сети 192.168.0.0/23 на
маршрутизаторе R1.
\item Настроить протокол RIP и включить анонсирование сети 10.10.0.0/23 на
маршрутизаторе R2.
\item Настроить протокол RIP и включить анонсирование сети 172.16.0.0/23 на
маршрутизаторе R3.
\item Указать интерфейсы Gi 1/0/3 маршрутизатора R1, Gi 1/0/3 маршрутизатора
R2, Gi 1/0/3 маршрутизатора R3, пассивными.
\item Проверить таблицу маршрутизации
\item Проверить сетевую связность VPC1, VPC2 и VPC1, VPC3 командами ping и
tracert.
\item Отключить соединение между маршрутизаторами R1 и R2 и проверить
сетевую связность между VPC1 и VPC2 ping и tracert.
\end{enumerate}

\begin{figure}[H] 
    \centering
    \includegraphics[width=0.75\textwidth]{top.png}
    \caption{Пример собранной сети}
\end{figure}

\begin{figure}[H] 
    \centering
    \includegraphics[width=0.75\textwidth]{table.png}
    \caption{Таблица с настройками}
\end{figure}

\endinput